% ============================================================
% ============================================================
% ONERA M6 -- OVERALL ASSESSMENT
% ============================================================
% ============================================================

\section[Overview]{\textbf{ONERA M6 Overall Assessment}}

% ============================================================
% ONERA M6 -- Overall Assessment
% ============================================================

\begin{frame}{\small ONERA M6 -- Overall Assessment}
\scriptsize

\begin{block}{\scriptsize Overall Assessment}
\begin{itemize}
    \item \textbf{Aerodynamics:} Limited grid and code-to-code variability overall, except for $C_D$.
    \item \textbf{Droplet Impingement:} Good overall agreement; impingement extent is more sensitive to spatial and droplet-distribution resolution than integrated water mass.
    \item \textbf{HTC \& Thermodynamics:} Substantial code-to-code variability, particularly in $HTC$, convective heat transfer, and freezing fraction.
    \item \textbf{Ice Accretion:} Comparable water impingement does not translate into comparable ice mass or final ice geometry.
\end{itemize}
\end{block}

\vspace{0.10cm}

\begin{alertblock}{\scriptsize Overall Takeaway}
Code-to-code variability increases through the icing simulation chain, with large differences emerging in thermodynamics and final ice accretion. HTC and thermodynamic modeling deserve closer investigation for these glaze conditions, while geometry-evolution methods also clearly influence the final ice shape, even when thermodynamic predictions are comparable.
\end{alertblock}

\end{frame}